\section*{Introduction générale}
\addcontentsline{toc}{section}{Introduction générale}
\label{sec:introduction_generale}

L'assurance-vie constitue le principal instrument d'épargne financière des ménages français. À la fin de l'année 2024, l'encours total atteignait 1989 milliards d'euros \cite{assurance_vie_chiffres_2025}, niveau record illustrant la place centrale de ce support dans le patrimoine des ménages. Dès janvier 2025, ce montant franchissait le seuil symbolique des 2000 milliards d'euros \cite{assurance_vie_chiffres_2025}, confirmant la progression continue de cette enveloppe.

Les flux d'alimentation de l'assurance-vie témoignent également de son dynamisme : en 2024, les versements bruts se sont élevés à 173 milliards d'euros, générant une collecte nette de 29,4 milliards d'euros \cite{assurance_vie_chiffres_2025}. Le seul mois de janvier 2025 a enregistré 17,3 milliards d'euros de versements, pour une collecte nette de 4,5 milliards, soit un maximum historique pour cette période \cite{assurance_vie_chiffres_2025}.

Instrument d'épargne à la fois liquide et adaptable, l'assurance-vie combine une fonction de capitalisation (fonds en euros, unités de compte) et une fonction de transmission patrimoniale. Ce double rôle lui confère une place singulière dans le système financier français et s'inscrit dans un contexte plus large de diversification de l'épargne des ménages, documenté par l'Observatoire des produits d'épargne financière de la Banque de France \cite{rapport_observatoire_epargne_2025}.

L'objectif du présent document est d'analyser, dans une perspective comparative, les régimes fiscaux applicables à l'assurance-vie et au compte-titres ordinaire (CTO). Bien que l'assurance-vie bénéficie d'un cadre fiscal spécifique, notamment en matière successorale (Code général des impôts, article 990 I 2025; Code général des impôts, article 757 B 2025), il convient de souligner que, dans de nombreux scénarios, le CTO présente un avantage relatif plus marqué. En particulier, l'absence de prélèvements sociaux au décès et la neutralisation des plus-values latentes confèrent au CTO une attractivité qui, selon la configuration patrimoniale et la durée de détention, peut excéder celle de l'assurance-vie \cite{lechiquier_2019}.

Il convient de rappeler que les modèles mobilisés ici reposent sur un ensemble d'hypothèses simplificatrices. D'une part, nous supposons que la fiscalité du CTO ne s'applique qu'au décès, en négligeant la matérialisation éventuelle de plus-values si des ventes intermédiaires étaient réalisées, ce qui confère un biais favorable au CTO dans notre cadre. D'autre part, les simulations considèrent des paramètres fiscaux figés (taux de prélèvements sociaux, barèmes successoraux, abattements), alors que ceux-ci sont susceptibles d'évoluer dans le temps. De plus, les frais de gestion sont traités comme constants et homogènes, alors qu'ils varient selon les contrats et selon l'allocation entre fonds en euros et unités de compte. Nous faisons également abstraction d'éléments comportementaux (arbitrages intra-contrat, réallocation des actifs, stratégie de liquidité). Les simulations reposent sur des trajectoires déterministes de rendement, sans intégrer la volatilité ni le risque de marché, ce qui limite leur portée interprétative à des comparaisons de long terme. Finalement, nous supposerons pour l'ensemble des simulations que les primes versées ont été effectuées strictement avant 70 ans, pour simplifier les interprétations, nous ne considèrerons pas le cas où il y des primes versées avant puis après cette barrière.

Afin de répondre à ces enjeux, le rapport est organisé de manière progressive.
